% =====================================================================================================
% Chương 2 — bản viết lại, phiên bản 5 (2026-10-13). Dùng như \input{noi-dung/06-chuong2} của main.tex.
% Các bản 1–4 giữ nguyên. Bản 5 xử lý bản kiểm nguồn Codex R35 (156 mục) và nhận xét R30 của người đọc, giảng viên
% (B1–B5, N1–N5); lý do từng thay đổi: reviews/CH2-REVISION-LOG.md; ghi chú: 06-chuong2-v5-NOTES.md.
% Công thức (2.1) tách dòng bằng aligned (bản dựng thử 13/10 báo tràn 80,5 pt); Bảng 2.1 hẹp lại 0,2 cm (tràn 1,9 pt).
% Mỗi dòng "% EV:" trỏ tới EVIDENCE-REGISTER.md (A01…K52), các phụ lục 1–2 (N-xx, X3-xx),
% EVIDENCE-REGISTER-ADDENDUM-5-CH2.md (C2-01…C2-24) hoặc trang của bài ghi tại chỗ ("tr." = trang in trên bài).
% R35 = codex-work/r35-ch2-check/R35-CH2-FACTCHECK-CODEX-v2.md (mã mục F, T, M, P, E; đề xuất C1–C14).
% =====================================================================================================

\chapter{Cơ sở lý thuyết và các công trình liên quan}
\label{chuong:coso}

Chương này trình bày bài toán ghép chuyến động mà khóa luận dựa vào, các cách nghiên cứu trước đã đo và ràng buộc
lời hứa với khách, và những gì đã biết về dự báo thời gian đi và giá trị của độ tin cậy. Lời hứa ở đây là giờ đón,
giờ trả báo cho khách; khóa luận tập trung vào lời hứa giờ trả (promised drop-off time). Luật C của khóa luận giới
hạn tổng độ lớn các lần bộ điều phối dời giờ trả của mỗi khách qua cả chuyến bằng một trần; định nghĩa hình thức của
nó và của các luật đối chứng nằm ở Chương~\ref{chuong:phuongphap}.
% EV: A06; K43; 07-chuong3-v4-v3.tex:18-25 (định nghĩa một câu của luật C ở Chương 3)

% -----------------------------------------------------------------------------------------------------
\section{Bài toán ghép chuyến động và kiến trúc yêu cầu--chuyến--xe}
\label{muc:c2-baitoan}

Trong dịch vụ đi chung xe, mỗi yêu cầu gồm điểm đón, điểm trả và thời điểm đặt; mỗi xe trong đội chở được nhiều
khách cùng lúc, nhưng không quá số chỗ. Bộ điều phối (dispatcher) quyết định nhận hay từ chối từng yêu cầu, giao nó
cho xe nào và xếp các điểm dừng theo thứ tự nào. Đây là một dạng của bài toán đưa đón theo yêu cầu (Dial-a-Ride
Problem, DARP): tìm các tuyến xe có tổng chi phí nhỏ nhất để chở chung khách theo đặt chỗ, đồng thời giữ các ràng
buộc về mức dịch vụ~\cite{molenbruch2017}. Ở dạng động, bộ điều phối chỉ biết một yêu cầu khi khách gửi
nó~\cite{schulz2026}, nên giải lại bài toán ở mỗi lượt quyết định (decision epoch) với trạng thái của đội xe lúc đó.
% EV: J25; C2-12 (Molenbruch, PDF tr. 3: định nghĩa DARP); C2-07 (Schulz–Pfeiffer tr. 549: yêu cầu không biết trước); X3-01; R35 F01–F05

Alonso-Mora và cộng sự~\cite{alonsomora2017} đề xuất một cách giải theo lô cho dịch vụ đi chung xe. Yêu cầu được gom
trong một khoảng ngắn, 30 giây trong thí nghiệm của họ, rồi được gán cùng một lần qua hai bước. Bước thứ nhất dựng
đồ thị yêu cầu--chuyến--xe (request--trip--vehicle graph), trong đó một chuyến là một nhóm yêu cầu mà một xe chở chung
được trong các ràng buộc dịch vụ; các nhóm được dựng dần từ nhỏ đến lớn. Hörl và Zwick~\cite{horl2022} và bộ mô
phỏng FleetPy~\cite{engelhardt2026} đều dùng thuật toán này.
% EV: J17; C2-10 (Alonso-Mora, bản PNAS sớm tr. 3: đồ thị yêu cầu–xe, yêu cầu–chuyến–xe, kiểm nhóm con; tr. 4–5: lô 30 s); C2-02 (Hörl–Zwick tr. 4, PDF tr. 5; tr. 9, PDF tr. 10: cài trong MATSim); J18; C2-11 (FleetPy tr. 11: thuật toán Alonso-Mora); R35 F06–F10

Bước thứ hai, bước chọn, là một quy hoạch nguyên (integer linear program). Gọi $\epsilon_{ij} \in \{0, 1\}$ bằng 1
khi xe $j$ được giao chuyến $T_i$, và $\chi_k \in \{0, 1\}$ bằng 1 khi yêu cầu $k$ không được xe nào phục vụ. Bài
toán là
\begin{equation}
\begin{aligned}
\min\quad & \sum_{(i,j)} c_{ij}\,\epsilon_{ij} + \sum_{k} c_{ko}\,\chi_k\\
\text{thỏa}\quad & \sum_{i} \epsilon_{ij} \le 1 \quad \text{với mọi xe } j,\\
& \sum_{(i,j):\, k \in T_i} \epsilon_{ij} + \chi_k = 1 \quad \text{với mọi yêu cầu } k .
\end{aligned}
\label{eq:c2-chon}
\end{equation}
Trong đó, tổng thứ nhất lấy trên các cặp chuyến--xe có trong đồ thị; $c_{ij}$ là tổng độ trễ của mọi khách khi xe $j$
chạy chuyến $T_i$, kể cả khách đang ngồi trên xe; $c_{ko}$ là một hằng số lớn phạt mỗi yêu cầu không được phục vụ (ký
hiệu theo bài gốc). Độ trễ của một khách là hiệu giữa giờ trả và giờ sớm nhất khách có thể tới nếu đi thẳng ngay lúc
gọi, nên gồm cả thời gian chờ. Ràng buộc thứ nhất cho mỗi xe nhận nhiều nhất một chuyến; ràng buộc thứ hai buộc mỗi
yêu cầu hoặc thuộc đúng một chuyến được chọn, hoặc không được phục vụ.
% EV: C2-10 (Alonso-Mora tr. 3: hàm chi phí (1), hằng số c_ko cho tập R_ko, định nghĩa độ trễ; tr. 4: giải thuật 1, dòng 2–4); ký hiệu c_ko kiểm lại trên PDF tr. 3 bằng pdftotext ngày 13/10; R35 F11–F15

Hàm mục tiêu~\eqref{eq:c2-chon} có tính độ trễ của khách đang trên xe, nhưng chỉ như một khoản chi phí cộng chung cho
mọi khách và đo so với chuyến đi thẳng, không so với giờ đã báo; ngoài độ trễ tối đa, bài không đặt trần riêng lên
tổng độ lớn các lần dời giờ trả của từng khách. Bài giữ giờ đón đã báo bằng một ràng buộc (mục~\ref{muc:c2-rangbuoc});
việc các giờ đã báo bị thay đổi thì Fielbaum và Alonso-Mora đo và kiểm soát (mục~\ref{muc:c2-dolech}).
% EV: C2-10 (tr. 3: chi phí (1) gồm độ trễ của khách trên xe, đo so với đi thẳng; tr. 4: giờ đón muộn nhất hạ về giờ đón dự kiến, khách trên xe luôn được trả trong độ trễ tối đa; công thức (1) và giải thuật 1 không có đại lượng nào cộng các lần dời); K30; R35 F16–F19, C5 (bỏ vế "miễn là mỗi lần tổng chi phí nhỏ hơn" của v4)

RideBound giữ bước chọn theo lối này nhưng không dựng đồ thị yêu cầu--chuyến--xe. Với mỗi xe, bộ sinh chỉ chèn thêm
(insertion-only generator) tạo các phương án: giữ nguyên kế hoạch cũ, hoặc chèn điểm đón, điểm trả của yêu cầu mới
vào kế hoạch cũ mà không bỏ hay đảo điểm cũ. Bộ chọn (selector) giao mỗi xe đúng một phương án và so các tiêu chí
theo thứ tự, trong đó tổng chi phí kế hoạch chỉ là một tiêu chí xếp sau (mục~\ref{muc:baitoan-phatbieu}); yêu cầu đã
được nhận không bị chuyển sang xe khác. Bảng~\ref{bang:c2-dichvu} so các ràng buộc dịch vụ của tập thử với các nghiên
cứu trên.
% EV: X3-03 (GEN:196-212, 231-246, 548-575); X3-22 (ORT:96-125); X3-07 (tổng chi phí kế hoạch là tiêu chí 4); X3-39 (không cấu hình nào chuyển khách đã nhận sang xe khác); A14; A24; 07-chuong3-v4-v3.tex:142-149 (Bảng 3.2); R35 F20–F22

\begin{table}[htbp]
\centering
\caption{Ràng buộc dịch vụ của một tuyến xe: cách đặt ở một số nghiên cứu trước và giá trị trong các tập thử của
khóa luận. Giá trị của khóa luận là cấu hình của tập thử, không phải chuẩn ngành; mỗi nghiên cứu tính từ một mốc
khác nên các con số không so trực tiếp được.}
\label{bang:c2-dichvu}
\small
\begin{tabular}{@{}p{2.4cm}p{3.4cm}p{5.7cm}p{3.1cm}@{}}
\toprule
\textbf{Ràng buộc} & \textbf{Ý nghĩa} & \textbf{Ở một số nghiên cứu trước} & \textbf{Trong các tập thử}\\
\midrule
Cửa sổ đón & Khách phải được đón trong một khoảng tối đa kể từ lúc gọi xe & 2, 5 hoặc 7
phút~\cite{alonsomora2017}; 5 phút~\cite{fielbaum2020,elmi2026}; 8 phút ở nghiên cứu điển hình của
FleetPy~\cite{engelhardt2026} & 10 phút kể từ lúc gọi\\
% EV: C2-10 (tr. 5: 120/300/420 s); C2-06 (FAM tr. 18, bảng 6); C2-08 (Elmi tr. 15); C2-11 (FleetPy tr. 11: 480 s); B03 (600 s); C2-22 (NORM:444); R35 T101–T105
Thời gian ngồi xe hay độ trễ & Giới hạn phần thời gian khách mất thêm so với đi thẳng & Độ trễ tổng, gồm cả thời gian
chờ, không quá hai lần cửa sổ đón~\cite{alonsomora2017}; độ trễ tổng tối đa 10 phút~\cite{fielbaum2020,elmi2026};
thời gian đi thẳng cộng một phần vượt tính bằng tỉ lệ của thời gian đi thẳng~\cite{molenbruch2017}; thời gian đi vòng thêm tối đa
40\% thời gian đi thẳng~\cite{engelhardt2026} & Thời gian ngồi trên xe không quá 1,5 lần thời gian đi thẳng\\
% EV: C2-10 (tr. 3, 5: δ ≤ ∆, ∆ = 2Ω, gồm thời gian chờ); C2-06 (bảng 6); C2-08 (tr. 15); C2-12 (PDF tr. 5: phần vượt là một tỉ lệ của thời gian đi thẳng); C2-11 (tr. 11: "maximum relative detour time of 40%"); B03; C2-22 (PhysicalPlanValidator.cs:594); R35 T106–T111, C7
Sức chứa & Số khách tối đa cùng lúc trên một xe & 1, 2, 4 hoặc 10 chỗ~\cite{alonsomora2017}; 4
chỗ~\cite{engelhardt2026}; 6 chỗ~\cite{schulz2026} & 4 chỗ\\
% EV: C2-10 (tr. 5); C2-11 (tr. 11: 180 xe 4 chỗ); C2-07 (tr. 562); B03; R35 T112–T115 (8 xe mỗi kịch bản ở Bảng 4.1, không phải sức chứa)
\bottomrule
\end{tabular}
\end{table}

Trong Bảng~\ref{bang:c2-dichvu}, độ trễ tổng của Alonso-Mora và cộng sự gồm cả thời gian chờ, còn giới hạn của khóa
luận chỉ tính thời gian ngồi trên xe; các giá trị của khóa luận, kể cả cửa sổ đón rộng hơn các giá trị đã trích, chưa
được hiệu chỉnh theo sở thích của khách.
% EV: C2-10 (tr. 5); C2-22 (PhysicalPlanValidator.cs:594: thời gian ngồi xe = giờ trả − giờ đón); B03; C2-23 (D03:252, 259-260); R35 F23, T116

% -----------------------------------------------------------------------------------------------------
\section{Đo lường và kiểm soát độ bất ổn của lời hứa}
\label{muc:c2-dolech}

Fielbaum và Alonso-Mora~\cite{fielbaum2020} chia độ bất ổn (unreliability) của dịch vụ đi chung xe theo nguồn: từ
khách, như khách vắng mặt hay hủy chuyến, và từ hãng xe, khi cách gán yêu cầu cho xe thay đổi mỗi lần có yêu cầu mới;
bài tập trung vào nguồn thứ hai. Với mỗi khách, giá trị báo ban đầu của thời gian chờ, độ vòng (thời gian ngồi xe trừ
thời gian đi thẳng) và độ trễ tổng là giá trị sẽ đạt nếu lộ trình của xe giữ nguyên từ lần gán thứ nhất. Chỉ số của
họ là phần giá trị thực tế vượt giá trị báo ban đầu, với giá trị âm đặt bằng 0 để chỉ giữ phần bất lợi; họ còn đếm số
lần giá trị dự kiến tăng lên.
% EV: J06; X3-35; C2-06 (FAM tr. 4: định nghĩa 1–2, tập trung vào hãng xe; tr. 5: giá trị báo ban đầu, độ vòng = thời gian trên xe − thời gian đường ngắn nhất, các hiệu, cắt giá trị âm; tr. 7: định nghĩa 9); R35 F24–F26

Thí nghiệm của họ dùng 19.610 yêu cầu taxi Manhattan trong khung 13--15 giờ ngày 15/01/2013, một khung thấp điểm được
chọn vì thời gian đi khá ổn định, để bỏ qua độ bất ổn do tắc đường. Khi không hạn chế việc chuyển khách sang xe khác,
hơn một phần ba số yêu cầu không bị từ chối ngay gặp ít nhất một thay đổi.
% EV: J07; C2-06 (tr. 9: 19.610 yêu cầu, 13–15 giờ ngày 15/01/2013, lý do chọn giờ thấp điểm; tr. 1, tóm tắt: "when reassignment is not restricted ... more than one third of the requests that are not immediately rejected face some change"); R35 F27–F29, C8 (giữ "hơn một phần ba" của tóm tắt, xem NOTES mục 6); F30 không còn trong văn bản

Ở mục 6, họ đưa việc kiểm soát vào bước gán xe. Mức mà bài gọi là kiểm soát hoàn toàn cấm từ chối khách đã gán và
hạ thời gian chờ tối đa của họ về giá trị dự kiến, nhưng không giữ độ vòng hay giờ trả. Đổi lại, theo bảng 5 của bài,
tỉ lệ yêu cầu bị từ chối tăng từ 34,5\% lên 45,6\%. Chi phí trung bình của khách, quy ra phút ngồi xe theo trọng số của
bài, cũng tăng từ 9,56 lên 11,6. Mức trung gian chỉ cấm từ chối khách đã gán: tỉ lệ từ chối còn 31,2\%, nhưng
thời gian chờ và độ vòng trung bình tăng khoảng 50\%. Bài để ngỏ câu hỏi kiểm soát độ bất ổn mà không tăng từ chối,
độ trễ hay thời gian chờ.
% EV: J08; J09; K32; C2-06 (tr. 17: hai mức kiểm soát; mức hoàn toàn chỉ cập nhật thời gian chờ tối đa, độ vòng trung bình vẫn tăng 0,65 → 0,94; bảng 5, "All times are in minutes": cơ sở 1,01/0,65/34,5%/9,56, trung gian 1,51/1,01/31,2%/10,3, hoàn toàn 0,77/0,94/45,6%/11,6, thứ tự hàng đã đối chiếu ảnh trang ở R35; "increased in about 50% on average"; tr. 18, bảng 6: trọng số một phút ngồi xe 1, một phút chờ 2, một lần từ chối 40; tr. 18: câu hỏi mở); R35 F31–F35, C8, P03

Elmi và cộng sự~\cite{elmi2026} tiếp tục hướng này theo từng khách. Độ chùng (slack) của một khách là khoảng còn lại
giữa giới hạn và giá trị dự kiến lúc gán, tính riêng cho thời gian chờ và cho độ trễ. Điểm khả năng đi chung của mỗi
khách ước lượng số yêu cầu có thể ghép vào chuyến của họ; khách có điểm thấp nhận độ chùng 0 nên không bị đổi giờ,
khách khác giữ toàn bộ độ chùng. So với cách siết đều cho mọi khách, ở tỉ lệ từ chối gần bằng nhau (29,1\% và 29,2\%,
trung bình trên 10 ngày), cách này bảo đảm ngay lúc gán cho 83,9\% số khách rằng giờ không đổi, còn cách siết đều gần
như không bảo đảm được cho ai, trong mô hình có thời gian đi trên mỗi đoạn đường không đổi.
% EV: C2-08; J12 (nâng ở ADD5); Elmi tr. 1 (tóm tắt); tr. 2, 5 (không xét giao thông; công thức 2; thời gian đi trên cạnh không đổi); tr. 11 (điểm tích phân từ lúc gán tới lúc trả = số yêu cầu có thể ghép); tr. 12–13 (cách đều, cách theo điểm 0% hay 100%; độ chùng 0 ngăn mọi thay đổi); tr. 16 (mỗi điểm là trung bình 10 ngày); tr. 18 (83,85% có độ chùng 0, tóm tắt làm tròn 83,9%; với cách đều "such a guarantee cannot be provided to anybody", 99,95% còn được phép đổi); tr. 19, bảng 2 (từ chối 29,09% và 29,21%; tóm tắt viết "the same rejection rate of 29%"); R35 F36–F41, C4; R30 N5

Thước đo của Fielbaum và Alonso-Mora so giá trị cuối với giá trị báo ban đầu, nên trong một ví dụ dựng tay, khách bị
dời muộn 5 phút rồi được kéo về giờ cũ có hiệu bằng 0; chỉ số đếm số lần tăng ghi lần dời đó nhưng không cho biết độ
lớn. Khóa luận kế thừa ý tách phần thay đổi do cách gán xe và kiểm soát nó trước khi quyết định, nhưng ở từng lượt so
giờ trả của phương án với giờ trả khi xe giữ nguyên kế hoạch, rồi cộng độ lớn phần chênh qua cả chuyến
(mục~\ref{muc:bamoc}).
% EV: C2-06 (tr. 5: hiệu giữa thực tế và báo ban đầu; tr. 7: chỉ đếm số lần tăng); ví dụ 5 phút dựng tay từ định nghĩa; A01; A04; X3-25; K30; K32; R35 F42–F43, C9. Câu "1 mili giây" của v4 bỏ khỏi Chương 2 (R30 B5, R35 F44): Chương 3 nêu nó với đúng phạm vi ở 07-chuong3-v4-v3.tex:297-300 (N-05; ADD1:104-111)

% -----------------------------------------------------------------------------------------------------
\section{Các dạng ràng buộc lời hứa}
\label{muc:c2-rangbuoc}

Ngoài cửa sổ đón và giới hạn thời gian ngồi xe, một số nghiên cứu đặt giới hạn trực tiếp lên giờ đã báo;
Bảng~\ref{bang:c2-dangrangbuoc} so các cách đó theo năm câu hỏi. Phần lớn chúng giữ cố định một giờ đã báo: giờ đón
dự kiến lúc ghép xe~\cite{alonsomora2017}, giờ đón hẹn ban đầu cộng một khoảng đệm~\cite{lu2024}, hay thời gian
chuyến đã báo lúc khách chấp nhận~\cite{laupichler2026}. Schulz và Pfeiffer~\cite{schulz2026} đặt hạn chót trả bằng
giờ trả trong lời mời (offer) cộng một độ chùng, mặc định 900 giây; với thời gian đi cho trước, họ thấy 10--15 phút là
đánh đổi hợp lý giữa tỉ lệ từ chối và sức hấp dẫn với khách đã nhận.
% EV: C2-10 (Alonso-Mora tr. 4); C2-14 (Lu–Kühnel tr. 4: công thức 3, giờ đón hẹn ban đầu cộng đệm); C2-09 (Mt-KaRRi tr. 5: thời gian chuyến so với thời gian chuyến đã chấp nhận); J10; C2-07 (Schulz–Pfeiffer tr. 549: thời gian đi cho trước; tr. 549–551: hạn chót; tr. 562: 900 s; tr. 569, mục 6.3.5 và Finding 5: quét 0–30 phút, l nhỏ làm tăng từ chối, l lớn làm dịch vụ kém hấp dẫn với khách đã nhận); J32; C2-09 (Mt-KaRRi tr. 1: bộ điều phối cho mô phỏng quy mô lớn; tr. 5; tr. 27: α = 1,4, β = 10 phút); R35 F45–F49, C2

Trong các cách ở Bảng~\ref{bang:c2-dangrangbuoc}, gần luật C nhất là cách Hörl và Zwick~\cite{horl2022} xử lý thời
gian đi thay đổi, thử trong bộ mô phỏng MATSim với nhu cầu lịch sử của dịch vụ MOIA ở Hamburg. Bài chỉ ra rằng khi
thời gian đi tăng, xe đang chở khách có thể không còn thứ tự trả khách nào kịp giờ trả muộn nhất, tức giờ tính từ thời
gian chờ tối đa và độ vòng cho phép. Vì vậy, ở đầu mỗi lượt họ tính lại giờ trả dọc tuyến đang gán bằng thời gian đi hiện
tại, gọi là giờ trả cập nhật. Theo công thức 7 của bài, với các yêu cầu đã gán cho xe đang xét, giờ trả bắt buộc là
giờ muộn hơn giữa giờ trả cập nhật và giờ trả muộn nhất; tuyến mới phải thỏa giờ này như ràng buộc cứng, còn phần vượt
giờ trả muộn nhất bị phạt theo từng giây (công thức 11). Như vậy, giờ trả chỉ được nới khi giờ trả cập nhật đã vượt
giờ trả muộn nhất; giờ đón được xử lý tương tự (công thức 8).
% EV: J01; J02; J03; J04; X3-36; C2-04; C2-05; Hörl–Zwick (trang in; PDF = trang in + 1), kiểm lại trên PDF 13/10 bằng pdftotext: tr. 4 dòng 34–35 (giờ muộn nhất từ chờ tối đa và độ vòng cho phép); tr. 7 dòng 8–17 (không còn thứ tự trả khả thi), dòng 31–40 (giờ trả cập nhật, phạt mỗi giây); tr. 8 dòng 8–21 (công thức 7 cho "all requests that are assigned to the currently evaluated vehicle", 8, 9–10, 11); tr. 10 (MOIA, MATSim); K26; R35 F50–F54, F84, C1, P01; R30 B1. Câu kết luận của bài chuyển sang mục 2.4

Ví dụ dựng tay: khách A được hứa trả lúc 8:30:00, giờ trả muộn nhất 8:35:00. Tắc đường làm giờ trả của A khi giữ
nguyên kế hoạch thành 8:32:00; ở hai lượt liên tiếp, hai lần ghép khách mỗi lần làm A tới muộn thêm 20 giây. Công thức
7 nhận cả hai lần vì A vẫn tới trước 8:35:00; luật C trần 30 giây loại lần thứ hai vì tổng phần dời thành 40 giây,
và không tính hai phút do tắc đường. Nếu giờ trả khi giữ nguyên kế hoạch đã là 8:36:00 thì ngược lại: công thức 7
không cho dời thêm, còn luật C vẫn cho dùng phần trần còn lại.
% EV: ví dụ dựng tay, áp đúng công thức 7 và 10 (Hörl–Zwick tr. 8) và định nghĩa luật C (A06; K43; 07-chuong3-v4-v3.tex:309-326); dựa trên ví dụ của R30 N1, thêm vế ngược lại cho công bằng

\begin{table}[htbp]
\centering
\caption{Các dạng ràng buộc lời hứa ở nghiên cứu trước và luật C, mô tả bằng lời. ``Không'' ở cột năm nghĩa là ràng
buộc so một giờ với một mốc nên các lần dời sớm, muộn có thể bù trừ nhau, dù một hạn chót cố định vẫn giới hạn phần
dời ròng. Cột ba là cách khóa luận đọc định nghĩa ràng buộc, không phải kết quả đo của các bài. Cửa sổ đón và giới
hạn ngồi xe ở Bảng~\ref{bang:c2-dichvu} không tách nguồn, nên phần do tắc đường bị tính.}
\label{bang:c2-dangrangbuoc}
\footnotesize
\setlength{\tabcolsep}{3pt}
\begin{tabular}{@{}p{2.6cm}p{2.7cm}p{2.9cm}p{2.3cm}p{1.8cm}p{2.6cm}@{}}
\toprule
\textbf{Dạng ràng buộc} & \textbf{So với mốc nào} & \textbf{Phần do tắc đường có bị tính?} & \textbf{Áp cho ai} &
\textbf{Cộng độ lớn các lần dời qua cả chuyến?} & \textbf{Dữ liệu giữ riêng cho từng khách}\\
\midrule
Giữ giờ đón đã báo~\cite{alonsomora2017,fielbaum2020,lu2024} & Giờ đón dự kiến lúc gán, hay giờ hẹn ban đầu cộng đệm
2--3 phút & Có, vì không tách nguồn; đệm chừa cho bất định & Khách chưa lên xe & Không & Giờ đón đã báo\\
% EV: C2-10 (tr. 4); C2-06 (tr. 17); C2-14 (tr. 4); R35 T207–T209
Hạn chót từ lời mời~\cite{schulz2026} & Giờ trả trong lời mời cộng độ chùng 900 giây & Ngoài phạm vi: thời gian đi
cho trước & Mọi khách đã nhận lời mời & Không & Hạn chót\\
% EV: J10; C2-07 (tr. 549–551, 562); R35 T210–T213, C2 (bỏ "khi thời gian đi đổi thì có" và dòng A10 của v4)
Hạn trả được nới (công thức 7)~\cite{horl2022} & Giờ trả muộn nhất, nới tới giờ trả cập nhật khi giờ này vượt nó &
Có, tới khi giờ cập nhật vượt hạn; phần vượt sau đó chỉ bị phạt & Các yêu cầu đã gán cho xe đang xét & Không & Giờ
trả muộn nhất, giờ trả cập nhật\\
% EV: J01; J03; C2-04 (tr. 7–8); R35 T214–T217, C1
Bảo đảm theo từng khách~\cite{elmi2026} & Thời gian chờ, độ trễ dự kiến lúc gán & Ngoài phạm vi: thời gian đi cố
định & Khách điểm thấp: độ chùng 0; khách khác: toàn bộ & Không & Điểm, độ chùng\\
% EV: C2-08 (tr. 2, 5, 11–13); R35 T218–T221
Giới hạn thời gian chuyến~\cite{laupichler2026} & 1,4 lần thời gian chuyến đã báo lúc chấp nhận, cộng 10 phút &
Ngoài phạm vi: thời gian đi cố định & Mọi khách đã chấp nhận & Không & Thời gian chuyến đã báo\\
% EV: J32; C2-09 (tr. 4, 5, 27); R35 T222–T225
\textbf{Luật C (khóa luận)} & Giờ trả nếu xe giữ nguyên kế hoạch, tính lại ở mỗi lượt & Không; phần dời do điều phối
vẫn tính với thời gian đi hiện tại & Mọi khách đã xác nhận đặt xe & Có & Phần trần đã dùng, ghi trong sổ\\
% EV: A01; A02; A06; A13; K43; X3-25; R35 T226–T228, C3
\bottomrule
\end{tabular}
\end{table}

Bảng~\ref{bang:c2-dangrangbuoc} cho thấy luật C dùng lại mốc giờ trả cập nhật của Hörl và Zwick nhưng giới hạn một
đại lượng khác: không đặt hạn cho giờ trả mà cộng độ lớn các lần bộ điều phối dời giờ trả so với mốc đó qua cả
chuyến, dưới cùng một trần cho mọi khách đã xác nhận đặt xe, và ghi phần đã dùng vào sổ của từng khách
(mục~\ref{muc:luatc}). Phần thay đổi khi giữ nguyên kế hoạch không bị tính, và trần không được nới khi đường đông lên.
Cộng có dấu cả phần dời do điều phối lẫn phần thay đổi khi giữ nguyên kế hoạch, để các lần dời sớm, muộn bù trừ nhau,
thì tổng bằng giờ trả mới trừ lời hứa đầu; giới hạn tổng này là luật hạn chót, chỉ mượn thành phần hạn chót của
Schulz và Pfeiffer với các mức trần như luật C thay cho 900 giây. Cộng độ lớn mọi thay đổi khách thấy thì được luật
trần độ dời khách thấy. Hai luật này làm đối chứng (mục~\ref{muc:doichung}). Bảo vệ từng khách, giới hạn giờ đến, chèn khách động và chuyển khách sang xe khác đều đã có
ở các nghiên cứu trên, nên khóa luận không nhận các ý đó là mới.
% EV: K01; J05; X3-36; A06; X3-25; K26 (không lấy cứng/mềm làm khác biệt); A12; 07-chuong3-v4-v3.tex:442-468 (eq:dongnhat, chuỗi liên tục v_j = p_{j−1}; tiền lệ Hörl–Zwick); J11 (mượn thành phần hạn chót, 30–60 s so với 900 s, không tái lập); D03:214-217; J38; J39; K46; K47; C2-07, C2-08, C2-09, C2-10; R35 F56–F59, C3, C9, C10, P02, P06; R30 B2, B3

% -----------------------------------------------------------------------------------------------------
\section{Dự báo thời gian đi và lời hứa theo dự báo}
\label{muc:c2-dubao}

Khi đường đông dần, giờ đón, giờ trả tính bằng thời gian đi lúc ra quyết định có thể sớm hơn giờ xe thực sự tới, dù
xe giữ nguyên kế hoạch. Hörl và Zwick xét đúng tình huống ước lượng thấp này và kết luận rằng khi đó hãng xe phải chọn
giữa từ chối khách đã nhận và trả muộn hơn giờ đã báo; với tốc độ thực đổi theo giờ, ảnh hưởng này nhỏ~\cite{horl2022}.
% EV: A05; F10; C2-05 (Hörl–Zwick tr. 11: sai số ước lượng −40% tới +20%; tr. 10: tốc độ GPS đổi theo giờ, cập nhật mỗi 60 phút; tr. 12 dòng 26–37: kịch bản tốc độ thực, độ trễ thêm trung bình dưới 1 giây; tr. 14 dòng 12–15: kết luận "effects are minor in a scenario with realistic travel times"); R35 F55, F60–F61, C11

Giờ đến dự kiến có vai trò quan trọng ở các nền tảng gọi xe~\cite{hu2022}. Uber dùng nó để tính giá cước, ước lượng
giờ đón, giờ trả và ghép khách với tài xế; mô hình học sâu DeeprETA của hãng hiệu chỉnh giờ đến do bộ tìm đường tính,
và các tác giả lưu ý giờ hiển thị cho khách có thể cần lấy theo một phân vị thay cho giá trị trung bình. Leonardi và
cộng sự~\cite{leonardi2024} dùng một số chuyến quan sát trong mô phỏng dịch vụ xe buýt theo nhu cầu để học ra một
khoảng thời gian chuyến, và đề xuất lấy mức dự báo được làm tiêu chí định quy mô dịch vụ. Vì vậy lời hứa theo dự báo
(forecast-aware promise), tức lời hứa đã cộng sẵn mức tăng tắc đường dự báo, không phải ý mới của khóa luận.
% EV: J35 (nguồn sơ cấp C2-15); C2-15 (DeeprETA tr. 1: "ride-hailing platforms", công dụng ở Uber, phân vị; tr. 3: bảng 1); J21; C2-13 (Leonardi PDF tr. 2: tóm tắt, xe buýt theo nhu cầu, mô phỏng VISUM ở hai thành phố Ý; tr. 4: định cỡ theo mức dự báo được); K37; R35 F62–F65, C11 (bỏ "hơn 90%")

Bộ dữ liệu chuẩn Manhattan của FleetPy~\cite{engelhardt2026} lấy nhu cầu từ chuyến taxi tuần 11--18/11/2018 và hiệu
chỉnh thời gian đi bằng một hệ số nhân cho mỗi giờ trong ngày, tính từ thời lượng chuyến ghi trong dữ liệu; gói dữ
liệu còn có tệp hệ số theo khung 15 phút cho từng ngày. Khóa luận dùng tệp này làm hệ số tắc đường (travel-time
factor) chung cho cả thành phố, không phải tốc độ từng đoạn đường (mục~\ref{muc:c4-kichban}). Vì cả thành phố dùng
một hệ số và mô phỏng không có nhiễu ngẫu nhiên, thời gian đi trong mô phỏng có thể dễ dự báo hơn ngoài đời; kết quả
của lời hứa theo dự báo ở Chương~\ref{chuong:thucnghiem} vì vậy có thể lạc quan hơn khi áp dụng thật.
% EV: J18; C2-11 (FleetPy tr. 10: tuần dữ liệu, hệ số theo giờ; tệp *_tt_factors.csv 96 dòng cách 900 s; PREF:51-55, 67-93); B06; I02 (bỏ chữ "mức trần" của I02); F10 (thế giới gần như tất định); hai vế "có thể" là suy luận của người viết; R35 F69–F73, C11; R30 B4

Dự báo cho lời hứa trong khóa luận chỉ dựa vào thông tin đã có lúc hứa: hồ sơ tắc đường, tức hệ số tắc đường theo
khung 15 phút của các ngày đã qua, và mức tắc đường đo được lúc đó (mục~\ref{muc:dubao}). Theo trang dữ liệu của Ủy
ban Taxi và Xe limousine New York (Taxi \& Limousine Commission, TLC), dữ liệu chuyến công khai hiện được công bố theo
tháng, thường trễ khoảng hai tháng~\cite{tlc}. Vì vậy, muốn lập hồ sơ từ các ngày ngay trước khi vận hành thật, hãng
xe cần nguồn dữ liệu kịp thời hơn, ví dụ dữ liệu chuyến của chính mình. Khóa luận giả định có nguồn như vậy và chưa
kiểm điều này, vì các kết quả chính chỉ tính lại lời hứa trên các lần chạy mô phỏng đã ghi.
% EV: F02; X3-11; 07-chuong3-v4-v3.tex:517-519 (g = 1 ở khung hiện tại), 547-552 (ngày nguồn), 554-556 (phát lại); B27; C2-16 (trang TLC tải 12/10/2026: "published monthly ... typically with a two-month delay"); F03 (chỉ phần phát lại); câu về dữ liệu của hãng xe là giả định, ghi rõ trong câu; R35 F66–F68, C6, P08; R30 N2

% -----------------------------------------------------------------------------------------------------
\section{Độ tin cậy có giá trị gì với khách}
\label{muc:c2-giatri}

Theo phân tích tổng hợp của Wardman và cộng sự~\cite{wardman2026}, một phút tới đích muộn so với giờ biểu hay giờ dự
kiến được coi nặng bằng trung bình 2,53 phút ngồi xe khi câu hỏi nêu một phân bố nhiều mức trễ (15 giá trị, 2
nghiên cứu), và 3,12 phút khi câu hỏi nêu một mức trễ cho một phần số chuyến (66 giá trị, 27 nghiên cứu). Với
dịch vụ đi chung xe, Alonso-González và cộng sự~\cite{alonsogonzalez2020} khảo sát 1.006 người ở Hà Lan, đo độ tin cậy
qua độ dao động của thời gian và qua khoản tăng thêm có hệ thống so với giờ được báo, và thấy tỉ số giữa giá trị độ
tin cậy và giá trị thời gian (reliability ratio) khoảng 0,5. Khoản tăng thêm nhỏ, không báo trước nhưng có hệ thống,
được định giá thấp hơn thời gian đã báo; người trả lời không được đặt vào tình huống tới muộn. Theo phần tóm tắt, trong
một thí nghiệm ngẫu nhiên trên 1.425.745 chuyến, Yu và cộng sự~\cite{yu2022} thấy cả độ lớn của thời gian chờ báo lúc
đầu lẫn tần suất cập nhật đều ảnh hưởng tới việc khách bỏ yêu cầu.
% EV: J23; C2-17 (Wardman tr. 2: hệ số so với thời gian ngồi xe; tr. 3, mục 14–15: MeanLate-Distn là phân bố 5–10 mức trễ, MeanLate-X/Y là một mức trễ cho một phần số chuyến; tr. 10, bảng 5: 2,53 (15; 2) và 3,12 (66; 27)); J24; C2-18 (Alonso-González tr. 1: N = 1.006, tỉ số khoảng 0,5; tr. 11: độ lệch chuẩn và phần dời có hệ thống, "unannounced yet systematic"; tr. 13: không có giờ đến mong muốn); J33; C2-19 (Yu–Zhang–Zhou: chỉ tóm tắt qua Crossref cr13: "magnitude of the initial WTI and the update frequency", "a major ride-sharing platform"); R35 F74–F79, C12

Ba kết quả này đo ba đại lượng khác nhau: tới muộn so với giờ biểu, độ dao động của thời gian, và thông tin thời gian
chờ trước khi lên xe. Chúng không cho giá tiền hay ngưỡng cảm nhận của trần 30 giây trên phần dời giờ trả do điều
phối, và khóa luận chưa có căn cứ bên ngoài để coi mức trần này là chuẩn dịch vụ; mỗi mức trần vì vậy được trình bày
như một điểm trên đường đánh đổi giữa số khách được phục vụ và độ ổn định của lời hứa.
Chương~\ref{chuong:phuongphap} đặc tả luật C trên nền các khái niệm này.
% EV: D20 (chỉ phần "chưa có căn cứ WTP/ROI bên ngoài"; bỏ câu "không có chuẩn ngành" theo R35 F82 và R30 N2); K52; J23 và J24 ("hai cấu trúc khác nhau"); C2-23 (D03:260); R35 F80–F83, C12
